% ============================================================
% Chapter: Mooring Design and Analysis
% Assessed MLOs: MLO-2
% Excellent-grade rubric criteria:
% - Mooring design (5%): the mooring system is clearly defined, analytically
%   consistent, and its neutral state, pretension, and load transfer are evidenced.
% ============================================================

\section{Mooring Design and Analysis}
\pagecheck{2--3}

% Rubric Checklist: Mooring design (5%).
\instructionplaceholder{Mooring system definition}{Describe the selected mooring arrangement and why it suits the concept and site constraints. Define line types, anchor strategy, fairlead locations, pretension targets, and key geometric or material parameters. Include a plan view and at least one labelled schematic.}

\subsection{Modelling assumptions and setup}
\instructionplaceholder{Required evidence}{State the assumptions used in the mooring model, including line behaviour, seabed interaction, hydrodynamic loading assumptions, and coupling strategy with the floater or structure. Explain why the assumptions are adequate for the design stage and identify known limitations.}

\subsection{Neutral configuration and pretension}
\instructionplaceholder{Required evidence}{Report the neutral-state geometry and pretension distribution. Show how pretension was selected and check that the resulting restoring characteristics are consistent with the target operational envelope and serviceability requirements.}

\subsection{Response under representative environmental conditions}
\instructionplaceholder{Required evidence}{Present mooring response for representative load cases, including peak line tensions, minimum tension checks, and governing utilisation trends. Explain how the mooring system contributes to restoring stiffness, force transfer, and overall structural behaviour.}
